% Partie : set-theory
% Chapitre : spine
% Section : foundation

\documentclass[../../../include/open-logic-section]{subfiles}

\begin{document}

\olfileid[fr]{sth}{spine}{foundation}

\olsection{Fondation}

Nous avons \emph{presque} terminé, mais pas \emph{tout à fait} :
rien, dans $\ZFminus$, ne garantit que \emph{tout} ensemble
appartienne à un certain $V_\alpha$, c'est-à-dire que chaque
ensemble se forme à une étape.

En un sens mathématique assez simple, la présence éventuelle
d'ensembles hors de la hiérarchie ne nous \emph{importe} pas : s'il
en existe, nous pouvons les ignorer. Mais nous avons motivé notre
\emph{concept} d'ensemble par l'idée que chaque ensemble se forme à
une étape (voir \stageshier{} dans la \olref[z][story]{sec}). Nous
voulons donc exclure la possibilité d'ensembles extérieurs à la
hiérarchie. Il nous faut, pour cela, ajouter un axiome qui assure que
chaque ensemble se trouve quelque part dans la hiérarchie.

Puisque les $V_\alpha$ sont nos étapes, nous pourrions simplement
envisager d'ajouter l'énoncé suivant comme axiome :

\begin{defish}
\emph{Régularité.} $\forall A \exists \alpha\, A \subseteq V_\alpha$
\end{defish}

Ce serait une démarche parfaitement raisonnable. Toutefois, pour
des raisons qui seront expliquées dans la section suivante, nous
adopterons plutôt un autre axiome :

\begin{axiom}[Fondation]
$(\forall A \neq \emptyset)(\exists B \in A)A \cap B = \emptyset$.
\end{axiom}

Avec quelques efforts, nous pouvons montrer, dans $\ZFminus$, que
la fondation entraîne la régularité :
\begin{defn}
Pour chaque ensemble $A$, posons :
\begin{align*}
	\text{cl}_0(A) &= A,\\
	\text{cl}_{n+1}(A) &= \bigcup \text{cl}_n(A),\\
	\text{trcl}(A) &= \bigcup_{n < \omega} \text{cl}_{n}(A).
\end{align*}
Nous appelons $\text{trcl}(A)$ la \emph{clôture transitive} de $A$.
\end{defn}
\noindent
Le nom de «~clôture transitive~» est bien choisi :
\begin{prop}\ollabel{subsetoftrcl}
$A \subseteq \trcl{A}$ et $\trcl{A}$ est un ensemble transitif.
\end{prop}

\begin{proof}
Il est clair que $A = \text{cl}_0(A) \subseteq \text{trcl}(A)$.
Si $x \in b \in \trcl{A}$, alors $b \in \text{cl}_n(A)$ pour un certain
$n$, et donc
$x \in \text{cl}_{n+1}(A) \subseteq \text{trcl}(A)$.
\end{proof}

\begin{lem}\ollabel{lem:TransitiveWellFounded}
Si $A$ est un ensemble transitif, alors il existe un ordinal
$\alpha$ tel que $A \subseteq V_\alpha$.
\end{lem}

\begin{proof}
Rappelons la définition de «~$\supstrict(X)$~» donnée dans la
\olref[ordinals][opps]{defsupstrict}. Définissons deux ensembles :
\begin{align*}
	D  &= \Setabs{x \in A}{\forall \delta\ x \nsubseteq V_\delta}\\
	\alpha &= \supstrict\Setabs{\delta}{(\exists x \in A)
	(x \subseteq V_\delta \land (\forall \gamma \in \delta)x \nsubseteq V_\gamma)}
\end{align*}
Supposons $D = \emptyset$. Pour chaque $x \in A$, il existe alors un
$\delta$ tel que $x \subseteq V_\delta$ et, par le bon ordre des
ordinaux, tel que
$(\forall \gamma \in \delta)x \nsubseteq V_\gamma$. Ainsi,
$\delta \in \alpha$, d'où $x \in V_\alpha$ par le
\olref[spine][Valphabasic]{Valphabasicprops}. Par conséquent,
$A \subseteq V_{\alpha}$, comme voulu.

Il suffit donc de montrer que $D = \emptyset$. Par l'absurde,
supposons le contraire. Par fondation, il existe
$B \in D \subseteq A$ tel que $D \cap B = \emptyset$. Si $x \in B$,
alors $x \in A$ puisque $A$ est transitif ; et, comme $x \notin D$,
il existe un $\delta$ tel que $x \subseteq V_\delta$. Posons donc
\[
	\beta = \supstrict\Setabs{\delta}{(\exists x \in B)(x \subseteq V_\delta \land (\forall \gamma < \delta)x \nsubseteq V_\gamma)}.
\]
\footnote{Note de l'édition française : l'original écrit ici
«~$x \in b$~», avec un $b$ minuscule non défini dans cette
démonstration. La formule doit porter sur l'ensemble $B$ choisi par
fondation.}
Comme précédemment, $B \subseteq V_\beta$, ce qui contredit
$B \in D$.
\end{proof}

\begin{thm}\ollabel{zfentailsregularity}
La régularité est vraie.
\end{thm}

\begin{proof}
Fixons $A$. Par la \olref{subsetoftrcl}, nous avons
$A \subseteq \trcl{A}$, et $\trcl{A}$ est transitif. Il existe donc
un $\alpha$ tel que
$A \subseteq \trcl{A} \subseteq V_\alpha$ par le
\olref{lem:TransitiveWellFounded}.
\end{proof}

Ces résultats montrent que $\ZFminus$ démontre l'implication
$\emph{Fondation}\Rightarrow\emph{Régularité}$. Dans la
\olref[spine][rank]{zfminusregularityfoundation}, nous montrerons que
$\ZFminus$ démontre aussi
$\emph{Régularité}\Rightarrow\emph{Fondation}$. Ainsi, la fondation
et la régularité sont \emph{équivalentes} modulo $\ZFminus$. Dans
cette théorie, nous pouvons donc justifier la fondation en remarquant
qu'elle équivaut à la régularité. Quant à la régularité, elle se
justifie immédiatement à partir de \stageshier{}.

\end{document}
